
The widely held assumption that contrastive SSL amplifies spurious correlations due to augmentation invariance is not supported on Waterbirds. Supervised ERM encodes spurious background features significantly more strongly than contrastive MoCo-v3 (ratio diff = 0.0247, d = 5.68, p\_bonf < 0.0001). The paradigm effect is highly significant across both datasets (Waterbirds ANOVA F = 35.99, p = 2.42 $\times 10^{-7}$; CelebA F = 5.51, p = 0.009), but the paradigm ranking reverses between datasets, establishing that spurious attribute type interacts with pretraining objective. Background-replacement augmentation in SimCLR training produces a 23.1\% spurious/task ratio reduction (p = 1.26 $\times 10^{-6})$, confirming augmentation as a functional causal lever.

\textbf{Contributions:} (1) First controlled four-paradigm comparison of spurious feature encoding in frozen ResNet-50 representations across two benchmarks. (2) Supervised label correlation identified as the dominant driver --- ERM and DINO show equivalent spurious encoding despite different pretraining objectives; label-agnostic MoCo-v3 shows measurably lower encoding. (3) Spurious attribute type $\times$ pretraining objective interaction established via cross-dataset ranking reversal. (4) Background-replacement augmentation confirmed as a functional causal intervention for spurious encoding reduction.

\textbf{Future directions:} (i) Comparison with masked autoencoder (MAE, label-free reconstruction) to distinguish label-agnostic training from contrastive-specific augmentation effects; (ii) Downstream worst-group accuracy evaluation via DFR on MoCo-v3 vs ERM features; (iii) ViT backbone comparison; (iv) Additional spurious attribute types to characterize the paradigm $\times$ attribute-type interaction matrix.

